% Ad Atticum 13.19 (ad-atticum-13-19)
% Source: https://raw.githubusercontent.com/PerseusDL/canonical-latinLit/master/data/phi0474/phi057/phi0474.phi057.perseus-lat1.xml
% Fetched by scripts/fetch_latin.py

% Dateline (Perseus): Scr. in Arpinati prid. K. Quint. a. 709 (45) .

\ciceroLetterOpener{CICERO ATTICO salutem}

\ciceroSection{1} commodum discesserat Hilarus librarius iv Kal., cui dederam litteras ad te, quom venit tabellarius cum tuis litteris pridie datis; in quibus illud mihi gratissimum fuit quod Attica nostra rogat te ne tristis sis, quodque tu ἀκίνδυνα esse scribis.

\ciceroSection{2} Ligarianam, ut video, praeclare auctoritas tua commendavit. scripsit enim ad me Balbus et Oppius mirifice se probare ob eamque causam ad Caesarem eam se oratiunculam misisse. hoc igitur idem tu mihi antea scripseras.

\ciceroSection{3} in Varrone ista causa me non moveret ne viderer φιλένδοξοσ (sic enim constitueram neminem includere in dialogos eorum qui viverent); sed quia scribis et desiderari a Varrone et magni illum aestimare, eos confeci et absolvi nescio quam bene, sed ita accurate ut nihil posset supra, Academicam omnem quaestionem libris quattuor. in eis quae erant contra ἀκαταληψίαν praeclare conlecta ab Antiocho, Varroni dedi. ad ea ipse respondeo; tu es tertius in sermone nostro. si Cottam et Varronem fecissem inter se disputantis, ut a te proximis litteris admoneor, meum κωφὸν πρόσωπον esset.

\ciceroSection{4} hoc in antiquis personis suaviter fit, ut et Heraclides in multis et nos in vi de re publica libris fecimus. sunt etiam de oratore nostri tres mihi vehementer probati. in eis quoque eae personae sunt ut mihi tacendum fuerit. Crassus enim loquitur, Antonius, Catulus senex, C. Iulius frater Catuli, Cotta, Sulpicius. puero me hic sermo inducitur, ut nullae esse possent partes meae. quae autem his temporibus scripsi Ἀριστοτέλειον morem habent in quo ita sermo inducitur ceterorum ut penes ipsum sit principatus. ita confeci quinque libros περὶ ut Epicurea L. Torquato, Stoica M. Catoni, Περιπατητικὰ M. Pisoni darem. ἀζηλοτύπητον id fore putaram quod omnes illi decesserant. haec Academica, ut scis, cum Catulo, Lucullo, Hortensio contuleram. sane in personas non cadebant; erant enim λογικώτερα quam ut illi de iis somniasse umquam viderentur. itaque ut legi tuas de Varrone, tamquam ἕρμαιον adripui. aptius esse nihil potuit ad id philosophiae genus quo ille maxime mihi delectari videtur, †easque† partis ut non sim consecutus ut superior mea causa videatur. sunt enim vehementer πιθανὰ Antiochia; quae diligenter a me expressa acumen habent Antiochi, nitorem orationis nostrum si modo is est aliquis in nobis. sed tu dandosne putes hos libros Varroni etiam atque etiam videbis. mihi quaedam occurrunt; sed ea coram.
